\chapter{Verification, Fidelity, and Independent Checking}
\label{ch:10}

The preceding chapters of this part have treated generation, experimentation, and modelling as instruments whose products require checking, and the present chapter treats the checking itself. Verification is not a single activity. It includes the formal checking of a proof, the replication of an experiment, the review of a design against a code, the audit of a financial statement, the validation of a model against data, the inspection of a structure after construction, and the monitoring of a system in service. What these have in common is that each is meant to establish, by a procedure independent of the one that produced the claim, that the claim is correct to a stated standard. The independence is essential, and so is the standard, and a large part of the chapter is concerned with the ways in which each can quietly fail. The distinction between a model that has been verified and one that has been confirmed is old in the earth sciences~\cite{oreskes1994}, and current risk-management guidance for artificial intelligence treats validation and monitoring as continuing obligations and not as one-time events~\cite{nist2023}.

The question of independence is subtler than it first appears. Two procedures are independent in the relevant sense when the circumstances that cause one to give a wrong answer are not the circumstances that cause the other to do so. They may share nothing in the way of personnel, funding, or method, and still be dependent, if they rely on the same data, the same assumptions about the system, or the same software library. The history of engineering contains many failures in which a design was reviewed by several parties who each relied on the same standard calculation, and in which an error in the calculation passed all the reviews. An independent check should therefore be defined by the structure of its possible errors and not by the identity of the parties that perform it. The practical corollary is that the most informative checks are those that use a different method, as when a numerical result is compared with an analytic bound, or a simulation with a measurement, or an automated prediction with the judgment of an experienced practitioner, because a different method is more likely to fail differently. Where verification is performed by a system of the same kind as the one that generated the claim, the probability of shared error is high, and the check should be treated as screening.

The question of the standard is a question about what is being claimed. A verification procedure checks a claim against a model of reality, and the claim established is that the thing satisfies a property in the model. Whether the thing satisfies the property in reality depends on how closely the model represents reality in respect of that property. The difference is the \emph{fidelity gap}, and it is the place at which formal verification, impressive as it is within its domain, reaches its limit. A control program may be proved to keep a system within safe bounds under a stated model of the plant, and the proof is a mathematical fact; whether the plant behaves as the model says is an empirical question that the proof does not address. This is not an objection to formal methods, which eliminate a large class of errors, but a reminder that their guarantees are conditional on assumptions that have to be validated by other means. The conditional form of the guarantee should be stated in the documentation, so that a reader knows what has been proved and what has been assumed.

A third consideration is the trusted base. Every verification procedure rests on components that are not themselves verified by the procedure: the proof checker, the measuring instrument, the reference standard, the regulatory body, the auditors. These components constitute the trusted base of the verification, and the reliability of the conclusion cannot exceed the reliability of the trusted base. Two principles guide its design. The first is that the trusted base should be as small as possible and as simple as possible, so that it can itself be examined by many parties. The second is that the number of components in it should be counted, because the failure probabilities of the components accumulate, and a procedure with many individually reliable components may be less reliable than one with few. In mathematics the proof assistants with small kernels illustrate the first principle, since the kernel is a short program whose correctness can be studied in isolation. In the social domain, the equivalent of a small kernel is a clear rule of evidence, stated in advance and applied by a body that is accountable for its application.

Verification also has an economic dimension that is frequently overlooked. It is a form of labor, and like other forms of labor it requires skilled people, equipment, and time, and it competes for resources with the activities it is meant to check. When generation becomes cheap and verification remains costly, the pressure to reduce the cost of verification will be strong, and the techniques by which the cost is reduced include sampling, automation, delegation to the generating party, and the relaxation of the standard. Each of these has a legitimate form and an illegitimate one. Sampling is legitimate when the sampling plan is designed to give a stated protection against defects and illegitimate when the samples are chosen by the party being checked. Automation is legitimate when the automated checker is itself validated and illegitimate when it is trusted because it is convenient. Delegation to the generating party, which takes the form of self-certification, is legitimate only for low-stakes matters and in the presence of penalties for false certification. The book's insistence on independent checking is accordingly an insistence that these savings should not be taken in domains where the consequences of error are severe or irreversible, and that the funding of independent checking should be treated as a necessary expense of any program whose output will be relied upon.

Finally, verification does not end when a thing is deployed. Systems in service encounter conditions that were not anticipated in the tests, and their components age. Monitoring in service is therefore a continuing part of verification, and it is supported by the design of the system: a system that is instrumented to report on its own condition, whose logs are preserved and open to review, and whose failures are investigated and published, supplies the evidence by which earlier claims are corrected. Aviation provides the model in the form of its practice of investigating accidents and incidents, publishing the findings, and amending standards accordingly, and the practice has contributed more to the safety of air travel than any single technical innovation. The corresponding practice for programs of the kind this book describes is to require, as a condition of authorization, a commitment to monitoring and disclosure over the life of the program, with an independent body empowered to halt operations when the evidence warrants.

\section*{Formalization}

Two propositions are given, the first on the accumulation of errors in the trusted base and the second on the transfer of a guarantee across the fidelity gap.

Let the trusted base of a verification procedure consist of components $1,\dots,k$, and let $B_i$ be the event that component $i$ fails in a way that could produce a wrong verdict, with $\Pr(B_i)=\varepsilon_i$. Let $W$ be the event that the verdict is wrong.

\begin{proposition}
If $W\subseteq\bigcup_{i=1}^kB_i$ then $\Pr(W)\le\sum_{i=1}^k\varepsilon_i$, with no assumption on the dependence among the $B_i$. If the $B_i$ are independent, $\Pr(\bigcup_iB_i)=1-\prod_i(1-\varepsilon_i)\le\sum_i\varepsilon_i$, and the difference is at most $\tfrac12\bigl(\sum_i\varepsilon_i\bigr)^2$. Hence for a trusted base of $k$ components each with $\varepsilon_i\le\varepsilon$ the error bound grows at most linearly in $k$.
\end{proposition}

\begin{proof}
The first bound is Boole's inequality. For independent events, $\Pr(\bigcup B_i)=1-\prod(1-\varepsilon_i)$, and by induction $\prod_i(1-\varepsilon_i)\ge1-\sum_i\varepsilon_i$, which gives the inequality. The bound on the difference follows from the Bonferroni inequality, $\Pr(\bigcup B_i)\ge\sum_i\varepsilon_i-\sum_{i<j}\varepsilon_i\varepsilon_j$, and $\sum_{i<j}\varepsilon_i\varepsilon_j\le\tfrac12(\sum_i\varepsilon_i)^2$.
\end{proof}

For five components with $\varepsilon_i=10^{-3}$, the bound is $5\times10^{-3}$ and the exact probability under independence is $1-0.999^5\approx4.99\times10^{-3}$. The value of the proposition lies not in the arithmetic but in the design rule: to reduce the error of a verdict by a factor $r$ one must either reduce every term of the sum or remove components from the base, and the second is generally cheaper.

For the fidelity gap, let a system be characterized by a parameter $m\in\mathcal M$ and a design by $x$, and let $\rho(x,m)\in\R$ be a \emph{robustness margin} for a required property, with the property satisfied if and only if $\rho(x,m)>0$. Let $m^\ast$ be the true system and $\hat m$ the model, and suppose the margin is Lipschitz in the model: $|\rho(x,m)-\rho(x,m')|\le L\,d(m,m')$ for a metric $d$.

\begin{proposition}
If verification in the model establishes $\rho(x,\hat m)\ge\mu>0$ and the fidelity gap satisfies $d(\hat m,m^\ast)\le\eta$, then the property holds for the true system whenever $L\eta<\mu$, and the guaranteed margin in the true system is at least $\mu-L\eta$. If $L\eta\ge\mu$ the verification is silent about the true system.
\end{proposition}

\begin{proof}
$\rho(x,m^\ast)\ge\rho(x,\hat m)-L\,d(\hat m,m^\ast)\ge\mu-L\eta>0$.
\end{proof}

The proposition has two practical consequences. First, the quantity that determines whether a formal guarantee has practical force is the ratio $\mu/(L\eta)$, which can be increased by building in margin (raising $\mu$), by seeking designs that are insensitive to the model (lowering $L$), or by improving the model through measurement (lowering $\eta$). Second, it explains why optimization against a model is dangerous, as noted in Chapter 8: an optimizer reduces the margin $\mu$ in the direction in which the model is most favorable, so that the optimized design has a small $\mu$ and the proposition's condition is least likely to hold. The remedy is to optimize a worst-case margin over a set $\{m:d(m,\hat m)\le\eta\}$, which restores the guarantee at the price of a design that is less than optimal for the nominal model.
